\section{The sextic seed and its compact returns}\label{sec:deg12}

The first five upper launches supply four brackets. Two further enclosures at height eight permit a common reference to close after the outer coefficients have been chosen. The three entrance differences below will be used again at the exterior maximum.

\subsection{The rational sextic seed, Taylor-defect soundness, and degrees six and seven}\label{subsec:seed}
Define $p_a(x)=ax+\sum_{i=2}^6 a_i x^i$ for real $a\in J$ with the following exact rational data
\begin{equation}\label{eq:seed-coeffs}
  \begin{aligned}
    (a_2,a_3,a_4,a_5,a_6) &= (13.8510158,\ 0.8651558,\ -27.5087995,\ -1.7814695,\ 20.3082601),\\
    a_- &= -0.003525444747007079096354237286,\quad J=[a_-,a_+],\quad a_+=a_-+10^{-30},
  \end{aligned}
\end{equation}
and sample heights $(h_0,h_1,h_2,h_3,h_4)=(0.5,\,3.7,\,3.88,\,4.03,\,4.36)$.
The five coefficients of degrees two through six are the seven-decimal
roundings of the corresponding coefficients in the Palomar sextic
\cite[Theorem~1 and Appendix~A]{Palomar2026}.
The rounding places this family outside the $10^{-40}$ coefficient box
of that source. Here $h_0=1/2$ is shared, while the four other compact
heights and the shooting interval for $a$ are chosen for the present
construction. The enclosures below certify this family directly.

\begin{lemma}[Certified sextic half-return enclosures]\label{lem:seed}
For $F=p_a$, all seven strict inclusions in Table~\ref{tab:seed} hold at the stated endpoints of $J$. Consequently, for every $a\in J$, $\D(h_0),\dots,\D(h_4)$ have strict alternating signs $+,-,+,-,+$, while $B_+(0,8)>B_-(2,8)$ at $a=a_-$ and $B_+(2,8)<B_-(0,8)$ at $a=a_+$.
\end{lemma}

\begin{table}[ht]
\centering\small
\caption{Certified open intervals for half-return differences of $p_a$ in~\eqref{eq:seed-coeffs}, verified in \texttt{checks/\allowbreak verify\_\allowbreak seed.py}.}\label{tab:seed}
\smallskip
\begin{tabular}{@{}llll@{}}
\toprule
Quantity & Parameter & Certified open interval & Accepted Taylor steps ($+$ / $-$) \\
\midrule
$\D(0.5)$ & $a_+$ & $(2,3)\cdot 10^{-7}$ & $142 + 141$ \\
$\D(3.7)$ & $a_-$ & $(-3,-1)\cdot 10^{-24}$ & $399 + 375$ \\
$\D(3.88)$ & $a_+$ & $(3,4)\cdot 10^{-27}$ & $430 + 407$ \\
$\D(4.03)$ & $a_-$ & $(-3,-2)\cdot 10^{-28}$ & $444 + 414$ \\
$\D(4.36)$ & $a_+$ & $(6,8)\cdot 10^{-32}$ & $475 + 448$ \\
$B_+(0,8)-B_-(2,8)$ & $a_-$ & $(2,3)\cdot 10^{-32}$ & $545 + 710$ \\
$B_+(2,8)-B_-(0,8)$ & $a_+$ & $(-8,-6)\cdot 10^{-32}$ & $729 + 519$ \\
\bottomrule
\end{tabular}
\end{table}

\begin{proof}
We prove mathematical soundness of the exact integer/rational certificate algorithm implemented in \texttt{checks/verify\_seed.py}. For each reflected polynomial $G(x)=p_a(sx)$ ($s=\pm 1$, $a\in\{a_-,a_+\}$), write $G'(x)=G'(0)+xq(x)$ with $q(x)=\sum_{i=0}^4 q_i x^i$. Setting $u=q_3/(2q_4)$, $v=-23/50$, $w=q_2-q_4(u^2+2v)$, and $z=(q_1-2q_4 uv)/(2w)$ completes the square twice:
\[
  q(x) = q_4(x^2+ux+v)^2 + w(x+z)^2 + r,\qquad r = q_0 - q_4 v^2 - w z^2.
\]
Exact rational arithmetic verifies $q_4>0$, $w>0$, $r>1$ (hence $xq(x)\ge 0$ for all $x\ge 0$), and $|G'(0)|+\sum_{i=0}^4|q_i|100^{-(i+1)}<1$. Thus $G'(x)\ge -1$ throughout $[-1/100,\infty)$. Along any two points in $x\ge -1/100$, the symmetric part of the Jacobian of $X_G(x,y)=(y-G(x),-x)$ is $\operatorname{diag}(-G'(x),0)\preceq I_2$: for any proposed polynomial path $P(\tau)=(p(\tau),v(\tau))$ on normalized time $\tau\in[0,1]$ representing physical step size $h=1/N_{\mathrm{inv}}<1$ ($t=t_0+h\tau$), the Euclidean error $e(\tau)=|z(t_0+h\tau)-P(\tau)|$ satisfies $\dd(e^2)/\dd\tau=2h\langle z-P,X_G(z)-X_G(P)\rangle-2\langle z-P,P'-hX_G(P)\rangle\le 2he^2+2e|P'-hX_G(P)|$ as long as $z_x\ge -1/100$. With $\|H\|_1=\sum_k|[H]_k|$ denoting the exact $\ell^1$ sum of polynomial coefficients in $\tau\in[0,1]$, $|P'(\tau)-hX_G(P(\tau))|\le L:=\|p'-h(v-G(p))\|_1+\|v'+hp\|_1$. Integrating $\dd e/\dd\tau\le he+L$ with $e^h\le(1-h)^{-1}$ yields
\[
  \sup_{\tau\in[0,1]} e(\tau) \le e_{\mathrm{out}} := \frac{e_{\mathrm{in}}+L}{1-h},
\]
valid by a first-exit continuation argument whenever the entire tube $\big[\min_{[0,1]}p-e_{\mathrm{out}},\,\max_{[0,1]}p+e_{\mathrm{out}}\big]$ (bounded conservatively via positive/negative coefficient sums of $p$) lies strictly inside $(-1/100,\infty)$.

On each step, \texttt{verify\_seed.py} proposes degree-$36$ polynomials $(p,v)$ on the dyadic grid $2^{-256}$, forms the \emph{exact} degree-$216$ composition $G(p(\tau))$ in integer arithmetic (covering truncation and rounding simultaneously), rounds $e_{\mathrm{out}}$ upwards on $2^{-256}$, certifies departure ($\dot x>0$ on step~$0$ when $x_0=0$), certifies $x>0$ on all interior steps, and on the final step certifies $p(0)>e_{\mathrm{in}}$, $p(1)<-e_{\mathrm{out}}$, and $-\dot x\ge m_{\mathrm{sp}}>0$ across the entire $(x,y)$-tube. Bisecting $p(\tau_*)\approx 0$ ($180$ steps) bounds $|z_x(t_0+h\tau_*)|\le b:=|p(\tau_*)|+e_{\mathrm{out}}$, so the true zero $\tau_{\mathrm{hit}}$ of $z_x$ satisfies $h|\tau_{\mathrm{hit}}-\tau_*|\le b/m_{\mathrm{sp}}$ and $|z_y(t_0+h\tau_{\mathrm{hit}})-v(\tau_*)|\le e_{\mathrm{out}}+M_x b/m_{\mathrm{sp}}$ ($M_x=\sup|z_x|$ on the step). Each of the $14$ terminal magnitudes is enclosed to width $<10^{-40}$, proving Table~\ref{tab:seed}. Finally, since $\partial_a\D<0$ by Lemma~\ref{lem:rotational}, $\D(h_0),\D(h_2),\D(h_4)>0$ at $a_+$ implies positivity on all of $[a_-,a_+]$, and $\D(h_1),\D(h_3)<0$ at $a_-$ implies negativity on all of $[a_-,a_+]$.
\end{proof}

\begin{lemma}[Certified translated incoming enclosures for $p_a$]\label{lem:seed-incoming}
Put $\Delta=3/500$. For $F=p_a$ and every $a\in J$, the incoming differences $I_{+,p_a}(y)-I_{-,p_a}(y-\Delta)$ lie in the exact decimal open intervals $(-1.1\cdot 10^{-7},-5\cdot 10^{-8})$ at $y=4/5$, $(0.003,0.008)$ at $y=0$, and $(-0.007,-0.001)$ at $y=-20$ (hence have signs $-,+,-$).
\end{lemma}

\begin{proof}
For $x\ge 2$ and $s=\pm 1$, omitting the positive term $a_2 x^2>0$ gives $p_a(sx)/x^6\ge a_6-|a_5|/2+a_4/4-|a_3|/8-\max_J|a|/32>8$, so $p_a(sx)>8x^6$. On $x\ge 2$ and $y\le 1$,~\eqref{eq:spatial-ode} satisfies $0<\dd y/\dd x=x/(p_a(sx)-y)\le x/(8x^6-1)\le 1/(7x^5)$, so integrating inward from $+\infty$ to $x=2$ places the arrival ordinate on $x=2$ inside $[y-\frac{1}{448},y]$ ($\int_2^\infty\dd x/(7x^5)=1/(7\cdot 4\cdot 16)=1/448$). Since $|p_a'(0)|=|a|<2$ and $\lim_{x\to+\infty}p_a(sx)=+\infty$ for every $a\in J$, Lemma~\ref{lem:half-returns}(iii) guarantees that from $(q,y_2)=(2,y_2)$ for \emph{every} $y_2\in\R$ (including $y_2\le 0$ at the $y=0$ and $y=-20$ boxes) the forward trajectory stays in $x(\tau)>0$ until hitting $\Sigma_-$ transversely at $(0,-B_{s,a}(2,y_2))$. On $x=2$ with $y_2\le 1$, $\dot x=y_2-p_a(2s)<1-8\cdot 2^6<0$, so for $y_2<y_2'\le 1$ trajectory non-crossing in $x>0$ keeps the $(2,y_2')$ arc strictly inside the $(2,y_2)$ arc down to $\Sigma_-$ and yields $B_{s,a}(2,y_2')<B_{s,a}(2,y_2)$, while Lemma~\ref{lem:rotational} for $q=2>0$ and $y_2\in\R$ gives $\sgn\partial_a B_{s,a}(2,y_2)=-s$ across all of $J\times\R$. Hence the minimum and maximum of $I_{+,p_a}(y)-I_{-,p_a}(y-\Delta)$ over $a\in[a_-,a_+]$ are bounded strictly by evaluating $B_+(2,y_+)-B_-(2,y_-)$ at the two extreme corners $(a_+,y,y-\Delta-\frac{1}{448})$ and $(a_-,y-\frac{1}{448},y-\Delta)$. Enclosing these $12$ compact trajectories via the validated integrator of Lemma~\ref{lem:seed} (\texttt{checks/verify\_entrance.py}) yields the three intervals.
\end{proof}

\begin{proof}[Proof of Proposition~\ref{prop:low-deg} for $n=6$ and $n=7$]
Take $F=p_a$ ($a\in J$). Choose $c_{\min}\in(0,1)$ so small that $\sup_{x\in[0,c_{\min}],\,a\in J}|p_a(sx)|<4$. For any starting height $K\ge 9$, descending from $y=9$ to $y=8$ along~\eqref{eq:scalar-u} uses $p_a(sx)>8x^6$ on $x\ge 2$ ($u\ge 2$): while $u\ge 2$ and $8\le y\le 9$, $\dd u/\dd(9-y)\le 9-8(2u)^3<-4u^3$, which reaches $u\le 2$ ($x\le 2$) in height descent at most $\int_2^\infty\dd u/(4u^3)=1/32<1$; meanwhile on $0\le x\le c_{\min}$ and $y\ge 8$, $\dd u/\dd(-y)=y-p_a(sx)>4$ drives $x$ above $c_{\min}$ within height descent $c_{\min}^2/8<1/8$ and keeps $x\ge c_{\min}$ down to $y=8$. Hence both $X_\pm$-arcs from $(0,K)$ cross $y=8$ with $x_\pm(8)\in[c_{\min},2]$, so monotonicity in $q$ (Lemma~\ref{lem:half-returns}(iii)) and the last two rows of Table~\ref{tab:seed} give $\D(K)>B_+(0,8)-B_-(2,8)>0$ at $a_-$ and $\D(K)<B_+(2,8)-B_-(0,8)<0$ at $a_+$, yielding a unique $a_K\in(a_-,a_+)$ with $\D(K)=0$.

Now scale $z=x/R$, $U=z^2/2$, $t=y/R^6$ ($R\gg 1$) in $p_a(x)$ itself: $\Phi_s(z):=R^{-6}p_a(sRz)=E(z)+sR^{-1}B(z)+O(R^{-2})$ on bounded $z>0$, where $E(z)=a_6 z^6$ ($A(z)=6a_6 z^4$), $B(z)=a_5 z^5$, and $\eps_R:=R^{-10}$ governs $\eps_R\dd U/\dd(-t)=t-\Phi_s(\sqrt{2U})$. Put $Y=E(1)=a_6$ and $T=a_6+1$, and close $K=R^6 T$ at $a_R\in(a_-,a_+)$.

Write $\phi_s(U)=\Phi_s(\sqrt{2U})$, $t_0=R^{-2}$ and
$\ell=\kappa R^{-2/3}$ with $\kappa>0$ small. For a fixed large $M$,
$\phi_s(M)>T$ bounds both arcs above in $U$. Uniformly for $a\in J$,
\[
 \sup_{[0,2\ell]}|\phi_s|\le t_0/2,\qquad
 cR^{-4/3}\le\phi_s'\le C,\qquad |\phi_s''|\le C\quad(\ell\le U\le M).
\]
Indeed the leading slope is $6a_6z^4$ and the relative errors from each
lower power $i$ have size $O((Rz)^{i-6})$ for $z\ge\sqrt{2\ell}$.
A further derivative $z^{-1}\partial_z$ is bounded on this range.
The root $r_s(t)\in[2\ell,M]$ for $t_0\le t\le T$ is unique there and
$|r_s'|\le CR^{4/3}$. In descending $t$, a launch reaches $\ell$ within
$2\ell R^{-8}$ since its speed below $\ell$ is at least $R^8/2$.
Thereafter the difference from $r_s$ has damping
$R^{10}[\phi_s]_{U,r_s}\ge cR^{26/3}$ and forcing bounded by $CR^{4/3}$.
Variation of constants gives $|U-r_s|=O(R^{-22/3})$ after a descent
of $2R^{-3}$ from either fixed launch level.

Let $d(t)=U_T(t)-U_Y(t)>0$ and $g_s=w_s(R^6T)-w_s(R^6Y)$.
The longer arc is already near $r_s(Y)\to1/2$, so $d(Y)$ is bounded
above and away from zero. Over the initial interval of length $2R^{-3}$
the secant is bounded by separation from the origin or by a derivative
bound on a fixed positive interval. On the common tracked interval
$[\phi_s]_{U_T,U_Y}=\phi_s'(r_s(t))+O(R^{-22/3})$.
Below $t_0$ the barrier $x\ge c_{\min}$ holds down to $y=8$ and
the only singular term in $\phi_s'$ satisfies
$|a|R^{-5}/z\le c_{\min}^{-1}R^{-4}$. This interval costs $O(R^{-2})$.
At height eight the terminal magnitude as a function of $x^2/2$ has
derivative bounded above and away from zero by transversality and
compactness on $[c_{\min}^2/2,2]\times J$. Hence $g_s$ is $R^2d(8/R^6)$
times a factor between fixed positive constants. Integrating
$\eps_R(\log d)'=[\phi_s]_{U_T,U_Y}$ and changing from $t$ to
$z=\sqrt{2r_s(t)}$ gives, uniformly in $a$,
\[
 \eps_R\log g_s=-\int_{z_0}^{z_{s,Y}}\Phi_s'(z)^2\,\dd z/z+O(R^{-2}),
 \qquad z_0=\sqrt{2r_s(t_0)}\asymp R^{-1/3}.
\]
Expanding the upper root $z_{s,Y}=1-sR^{-1}B(1)/E'(1)+O(R^{-2})$ adds the moving-endpoint contribution $-sR^{-1}A(1)B(1)$ to $2sR^{-1}\int_0^1 A(z)B'(z)\,\dd z$ (the omitted leading even and linear-in-$R^{-1}$ integrals on $[0,z_0]$
are $O(z_0^{10}+R^{-1}z_0^9)$). Here the actual physical integral
starts at $z_0>0$. In its expansion the term from the square of the linear
coefficient costs $O(R^{-10}|\log z_0|)$ on $[z_0,2]$.
All other terms of order $R^{-2}$ or smaller have integrable powers of $z$
there. This gives $\eps_R\log\big(w_s(R^6 T)-w_s(R^6 Y)\big)=-\int_0^1 E'(z)^2\,\dd z/z - sR^{-1}W(1)+O(R^{-2})$, where $W(1)=2\int_0^1 AB'\,\dd z-A(1)B(1)=\int_0^1(AB'-A'B)\,\dd z=\frac23 a_6 a_5<0$ since $a_6>0>a_5$. With $f_e=A$ and $F_o=B$, this is
$W(1)=-I_1(1)$ in the notation of De Maesschalck and Huzak
\cite[Proposition~1]{DMH2015}. The cutoff estimate above gives the
$O(R^{-2})$ error uniformly in $a$.

At $a=a_R$, $\D(R^6 Y)=(w_-(R^6 T)-w_-(R^6 Y))-(w_+(R^6 T)-w_+(R^6 Y))<0$ for large fixed $R$ ($R^6 Y>4.36$). Applying Proposition~\ref{prop:sard}(i) at $(0.5,3.7,3.88,4.03,4.36,R^6 Y,R^6 T)$ with $\sigma_{\mathrm{det}}=+1$ (keeping $a_{\mathrm{reg}}\in(a_-,a_+)$) gives $6$ hyperbolic cycles on an open exact-degree-$6$ set ($\HH(6)\ge 6$), and applying Proposition~\ref{prop:sard}(ii) to the resulting positive-ending $7$-sample chain gives $7$ hyperbolic cycles on an open exact-degree-$7$ set ($\HH(7)\ge 7$).
\end{proof}

\section{A moving pair of folds over the sextic}\label{subsec:deg10}
In both the degree-ten construction of Proposition~\ref{prop:low-deg} and the degree-twelve construction of Theorem~\ref{thm:deg12}, a stationary inflection at $z=1$ ($u_0=1/2$) splits at width $\sqrt\eta$ while $L\asymp\eta^{-1}$ couples $\eps=L^2/S^2\asymp\eta^{10}$ directly to $\eta$, and the sextic seed $p_a$ contributes lower powers whose compact critical points sit at $z=O(\eta)$. We prove all shooting, cutoff, fold-localization, and gap estimates directly for the explicit profile template shared by both applications.

\begin{proposition}[Moving-layer upper gap expansion at coupled scale $\eps\asymp\eta^{10}$]\label{prop:moving-layer}
Let $\mathcal{P}$ be a compact parameter space, let $\kappa_S>0$ be constant, and let $\kappa_L(\pi)>0$ and $z_b(\pi)>1$ depend continuously on $\pi\in\mathcal{P}$. For $0<\eta\ll 1$, $\pi\in\mathcal{P}$, and $a\in J=[a_-,a_+]$ from~\eqref{eq:seed-coeffs}, set $L=\kappa_L(\pi)\eta^{-1}$, $S=\kappa_S L^6$, $\eps=L^2/S^2=\kappa_S^{-2}L^{-10}\asymp\eta^{10}$, and consider
\begin{equation}\label{eq:seeded-template}
  \begin{aligned}
    F_{\eta,a}(x;\pi) &= F_\eta^0(x;\pi) + ax,\\
    F_\eta^0(x;\pi)   &= S\Big(E_\eta(x/L;\pi) + \eta B(x/L;\pi) + \eta^{3/2}H(x/L;\pi)\Big) + a_2 x^2 + a_3 x^3 + a_4 x^4,
  \end{aligned}
\end{equation}
where $E_0(0;\pi)=0$, $E_\eta(z;\pi)=E_0(z;\pi)+\int_0^z v\big(\eta D_1(v;\pi)+\eta^{3/2}Z(v;\pi)\big)\,\dd v$, and the polynomials $A_0(z;\pi)=E_0'(z;\pi)/z$, $D_1,Z\in z^4\R[z^2]$, $B\in z^5\R[z^2]$, and $H\in z^7\R[z^2]$ have fixed degree bounds independent of $\pi$ and coefficients continuous
on $\mathcal{P}$. Assume the following conditions.
\begin{enumerate}
  \item[\textup{(a)}] \emph{Core-to-branch positivity and seed matching.} $A_0(z;\pi)=z^4(1-z^2)^2\widetilde Q(z^2;\pi)$ with $\widetilde Q(0;\pi)=6a_6/\kappa_S$ and $\min_{\pi\in\mathcal{P},\,z\in[0,z_b(\pi)]}\widetilde Q(z^2;\pi)=:c_Q>0$, while $\kappa_S\kappa_L(\pi)[z^5]B(\cdot;\pi)=a_5$.
  \item[\textup{(b)}] \emph{Inflection splitting and prefix moments.} With $K(\pi)=\frac43\widetilde Q(1;\pi)>0$, one has $B(1;\pi)=0$, $B'(1;\pi)=-K(\pi)$, $D_1(1;\pi)=-2K(\pi)$, $H(1;\pi)=8K(\pi)$, $\int_0^1 A_0 B'\,\dd z=\int_0^1 A_0 H'\,\dd z=\int_0^1 D_1 B'\,\dd z=0$, and $4\int_0^1\big(ZB'+D_1 H'\big)\,\dd z=K(\pi)^2\Lambda$ ($\Lambda=493/2$).
\end{enumerate}
Write $\Phi_s(z;\pi,a)=F_{\eta,a}(sLz;\pi)/S$, $\phi_s(u)=\Phi_s(\sqrt{2u})$, and $P_e=(\phi_++\phi_-)/2$, and let $T(\pi)\in\big(E_0(1;\pi)+\delta_1,\,E_0(z_b(\pi);\pi)-\delta_1\big)$ ($\delta_1>0$) be continuous on $\mathcal{P}$. Then there exist a fixed cutoff $C_0\ge 2$ and $\eta_0>0$ such that for all $\eta\in(0,\eta_0]$:
\begin{enumerate}
  \item[\textup{(i)}] Uniformly in $\pi\in\mathcal{P}$, $\D(ST(\pi),a_-;\pi)>0>\D(ST(\pi),a_+;\pi)$, and there exists a unique, jointly continuous map $\pi\mapsto a_\eta(\pi)\in\operatorname{int}J$ solving $\D(ST(\pi),a_\eta(\pi);\pi)=0$, across which $(h_0,\dots,h_4)$ retain strict signs $+,-,+,-,+$.
  \item[\textup{(ii)}] Uniformly in $(\pi,a)\in\mathcal{P}\times J$, $\sup_{0\le z\le C_0/L}|\Phi_s(z)|=O(\eta^6)\ll\eta^3$, and on the tracked high-branch interval $z\in[C_0/L,z_b(\pi)]$ (hence on $\{\Phi_s\ge\eta^3,\,z\le z_b(\pi)\}$) each side $\phi_s$ has \emph{only} two critical points, namely nondegenerate folds at heights $\widehat\mu_s<\widehat M_s$ near $P_e(1/2)$ converging in $u=1/2+\sqrt\eta\,v$ to those of $v^3+(-2-s)v+8s$ (whereas any compact critical points of $p_a$ lie strictly inside $0<z<C_0/L$ at scaled height $O(\eta^6)$). At $t=\eta^3$, the unique root $z_s(\eta^3)\in[C_0/L,1/2]$ satisfies $z_s(\eta^3)\asymp\sqrt\eta$.
  \item[\textup{(iii)}] Uniformly in $(\pi,a)\in\mathcal{P}\times J$, at the eight levels $t_m(\pi)=P_e(1/2;\pi)+K(\pi)\eta^{3/2}\xi_m$ of Table~\ref{tab:cubic}, the positive gaps $g_s^U=w_s(ST(\pi))-w_s(St_m(\pi))>0$ satisfy $-\eps\log g_s^U=\J_s^U(t_m(\pi))+O(\eta^3)$, where $\J_s^U$ integrates $(\phi_s'(u))^2\,\dd u$ from $z=C_0/L$ with the trapped-launch and jump deletions of~\eqref{eq:JU-def} at $\widehat\mu_s,\widehat M_s$. The calibration $B'(1)=-K$, $H(1)=8K$ interchanges the two sides of the cubic profiles in Table~\ref{tab:cubic}. Thus $\J_+^U(t_m)-\J_-^U(t_m)=-K(\pi)^2\eta^{5/2}D_{\mathrm{cub}}^U(\xi_m)+o(\eta^{5/2})$ uniformly on $\mathcal{P}$, so at $a=a_\eta(\pi)$ the mismatch $\D(St_m(\pi))$ alternates $-,+,-,+,-,+,-,+$ across $m=0,\dots,7$.
\end{enumerate}
\end{proposition}

\begin{proof}
(i)~By~\eqref{eq:seeded-template} and~(a), $S[z^6]E_0/L^6=a_6$, $S\eta[z^5]B/L^5=a_5$ ($[z^5]H=0$), $[x^6](F_\eta^0-SE_0(x/L))=O(\eta)$, and $[x^k]F_\eta^0=O(SL^{-k})=O(\eta^{k-6})$ for $k\ge 7$; thus $F_{\eta,a}(\cdot;\pi)\to p_a$ in $C^\infty$ on compact subsets of $\R$ uniformly on $\mathcal{P}\times J$. On $[0,z_b(\pi)]$, $E_0'(z;\pi)\ge c_Q z^5(1-z^2)^2$ gives a uniform constant $c_E>0$ with $E_0(z;\pi)\ge c_E z^6$ (using $(1-z^2)^2\ge 9/16$ on $[0,1/2]$ and $E_0(z)\ge E_0(1/2)$ on $[1/2,z_b(\pi)]$). For $x\in[X_0,Lz_b(\pi)]$ ($z=x/L$), dividing~\eqref{eq:seeded-template} by $x^6$ gives $SE_0(z)/x^6\ge\kappa_S c_E$ while the $D_1,Z$ integral contributes $O(\eta)$, the $B,H$ terms contribute $O\big((Lz)^{-1}+\eta^{3/2}z\big)=O(X_0^{-1}+\eta^{3/2})$ (as $L\eta=\kappa_L(\pi)$), and the degree-$\le 4$ seed terms contribute $O(X_0^{-2})$. Fixing $X_0\ge 2$ large and combining with $p_a(sx)>8x^6$ on $[2,X_0]$ yields a uniform $c_0>0$ such that $F_{\eta,a}(sx;\pi)\ge c_0 x^6$ on $[2,Lz_b(\pi)]$ for all $(\pi,a)\in\mathcal{P}\times J$ and small $\eta$.

Because $\Phi_s(z_b(\pi))>T(\pi)+\delta_1/2$ for small $\eta$, $x=Lz_b(\pi)$ confines forward arcs from $(0,ST(\pi))$. On $8\le y\le 9$, $F_{\eta,a}(sx;\pi)\ge c_0 x^6$ brings $x$ below $X_0$ in $O(X_0^{-4})$ height descent and $p_a(sx)>8x^6$ brings $x$ below $2$ in $\le 1/32$ further height descent, while a fixed $c_{\min}\in(0,1)$ with $\sup_{[0,c_{\min}]}|F_{\eta,a}|<4$ drives $x$ into $[c_{\min},2]$ at $y=8$ just as in the $n=6$ proof above. Uniform compact convergence to $p_a$ preserves the two strict shooting margins of Lemma~\ref{lem:seed} over $\mathcal P$. It also preserves the five compact signs on $\mathcal P\times J$, since their minimum absolute margin on $J$ is positive. Since $a$ enters $F_{\eta,a}=F_\eta^0+ax$ solely via $+ax$ and the transverse hits on $\Sigma_-$ depend continuously on $(\pi,a)$, Lemma~\ref{lem:rotational} makes the residual strictly decreasing in $a$. The face inequalities give one interior zero for each $\pi$. Compactness of $J$ and uniqueness identify the limit of the roots for every convergent parameter sequence (or net), giving a continuous $a_\eta:\mathcal{P}\to\operatorname{int}J$ solving $\D(ST(\pi),a_\eta(\pi);\pi)=0$.

(ii)~Differentiating $\Phi_s(z)$ on $z>0$ gives
\[
  \begin{aligned}
    \phi_s'(u) = \frac{\Phi_s'(z)}{z} &= A_0(z) + \eta D_1(z) + \eta^{3/2}Z(z) + s\eta\frac{B'(z)}{z} + s\eta^{3/2}\frac{H'(z)}{z}\\
    &\quad + \frac{4a_4 L^4 z^2 + 3sa_3 L^3 z + 2a_2 L^2 + saLz^{-1}}{\kappa_S L^6}.
  \end{aligned}
\]
Since $B\in z^5\R[z^2]$ and $H\in z^7\R[z^2]$, on $z\in(0,1/2]$ the ratio $|\phi_s'(u)-A_0(z)|/z^4$ is bounded by $C\big(\eta + \eta z^{-1} + (Lz)^{-2} + (Lz)^{-5}\big) \le C'\big(\eta + (Lz)^{-1}\big)$ uniformly on $\mathcal{P}\times J$, while $A_0(z)\ge\frac{9}{16}c_Q z^4$. Fix $C_0\ge 2$ large enough that $C'/C_0\le\frac{9}{64}c_Q$. On $[0,C_0/L]$ (physical $x=Lz\in[0,C_0]$), $|F_{\eta,a}(sx;\pi)|\le M(C_0)$ gives $\sup_{0\le z\le C_0/L}|\Phi_s(z)|\le M(C_0)/S=O(\eta^6)\ll\eta^3$, confining any compact folds of $p_a$ below $z=C_0/L$ and below $t=\frac12\eta^3$. On $[C_0/L,1/2]$ ($Lz\ge C_0$), we obtain $\frac{9}{32}c_Q z^4\le\phi_s'(u)\le C_{\mathrm{up}}z^4$; integrating $z\phi_s'(z^2/2)$ from $C_0/L$ gives $\Phi_s(z)\asymp z^6+O(\eta^6)$, so the unique root $z_s(\eta^3)$ at $t=\eta^3$ lies in $(C_0/L,1/2)$ with $c_{\mathrm{cut}}\sqrt\eta\le z_s(\eta^3)\le C_{\mathrm{cut}}\sqrt\eta$. On $[1/2,z_b(\pi)]$, $A_0(z)=4z^4\widetilde Q(z^2)(u-1/2)^2\ge\frac14 c_Q(u-1/2)^2$ while $|\phi_s'(u)-A_0(z)|=O(\eta)$; fixing $K_0$ large gives $\phi_s'(u)\ge\frac18 c_Q(u-1/2)^2>0$ outside $|u-1/2|\le K_0\sqrt\eta$. Inside $u=1/2+\sqrt\eta\,v$ ($|v|\le K_0$), Taylor expanding $E_0(\sqrt{2u})=E_0(1)+K(\pi)(u-1/2)^3+O((u-1/2)^4)$ and using~(b) (with the physical $a_4,a_2$ shifts absorbed in $P_e(1/2)$ plus $O(\eta^{5/2}v)$ and the $a_3,a$ odd terms bounded by $O(L^{-3})=O(\eta^3)$) shows that $(\phi_s(u)-P_e(1/2))/(K(\pi)\eta^{3/2})\to v^3+(-2-s)v+8s$ in $C^3$ uniformly on $\mathcal{P}\times J$. Hence on $[C_0/L,z_b(\pi)]$, each side $\phi_s$ has \emph{only} the two simple folds $\widehat\mu_s<\widehat M_s$ converging to those of $v^3+(-2-s)v+8s$. Moreover, differentiating $\phi_s'(u)$ once more in $u$ ($z^{-1}\partial_z$) produces smooth bounded polynomial terms plus $O(L^{-3}z^{-1}+L^{-5}z^{-3})=O(\eta^{5/2}+\eta^{7/2})$ on $z\ge\frac12 z_s(\eta^3)\asymp\sqrt\eta$, so $|\phi_s''(u)|\le C_1$ uniformly on $[z_s(\eta^3)^2/2,z_b(\pi)^2/2]$.

(iii)~Use descending height $\ell=-t$ as time, so that
$\eps U_\ell=t-\phi_s(U)$. All bounds below are uniform over
$\mathcal P\times J$ and the eight specified samples. Put
$h=\eta^3$ and $w=\eta^{7/2}$. Excise $t<h$ and the two bands of
radius $4h$ about the \emph{actual} critical heights. The error in
approximating a critical height by its cubic limit need not be
smaller than $h$.

\emph{Retained branches.}
At an actual fold $U_f$, $|\phi_s''(U_f)|\asymp\sqrt\eta$ by the
$C^3$ convergence in (ii). At a band edge the displacement of the
nearby root from $U_f$ is comparable to $\eta^{5/4}$, and its slope
has magnitude comparable to $\eta^{7/4}$. These estimates follow
by Taylor expansion about the actual critical point in the $v$ chart,
where the curvature is bounded above and away from zero. The regular
parts of that chart have slope bounded below in magnitude by $c\eta$.
Below the chart the innermost retained root has $z\gtrsim\sqrt\eta$
and $\phi_s'\gtrsim z^4$ near zero. Together with (ii) this gives,
on every retained rising branch $r(t)$,
\[
 \phi_s'(r)\ge c\eta^2,\qquad |r'(t)|\le C\eta^{-2}.
\]
Here $r$ denotes the root in $U$, with $r(h)\asymp\eta$ at the bottom.
The bound on $\phi_s''$ in (ii) also holds in the $w$-neighborhood
of these branches by its proof. For $e=U-r(t)$, at $|e|=w$,
\[
 \frac{\dd |e|}{\dd\ell}\le
      -c\eta^2w/\eps+C\eta^{-2}<0 .
\]
The restoring term has size at least $c'\eta^{-9/2}$. Thus the tube
stays above the compact cutoff and is invariant on each retained arc.

\emph{Capture from a launch.}
Each sample is a distance at least $c\eta^{3/2}$ from both critical
heights, by Table~\ref{tab:cubic}. Consider the first rising root
during the first $h$ of descent. Below $r(t)-w$ the residual
$t-\phi_s(U)$ is at least $c\eta^2 w$. On the adjacent monotone
piece this follows by integrating the slope over its last $w$ up to
the root. Any intervening local maximum lies below $t$ by at least
$c'\eta^{3/2}$, and the compact core has height $O(\eta^6)$.
Consequently the distance to the tube decreases at speed at least
$c\eta^{11/2}/\eps-C\eta^{-2}$, starting a bounded distance
away. Capture takes $O(\eps\eta^{-11/2})=O(\eta^{9/2})=o(h)$.
The same argument applies to the reference at $T$, whose initial
root is regular a fixed distance above the fold heights.

\emph{Passage through a moving minimum.}
Write $\mu=\widehat\mu_s$ and let $r_\mu$ be its inner landing in $U$.
For $\mu-4h\le t\le\mu-2h$ and $U\in[r_\mu,z_b^2/2]$ one has
$\phi_s(U)\ge\mu$, including at the old minimum. Thus an arc still
there at $t=\mu-2h$ obeys
\[
 U_\ell\le-2h/\eps .
\]
It reaches $r_\mu$ within $O(\eps/h)=O(\eta^7)$ of height.
Between $r(t)+w$ and $r_\mu$ the inner branch rises, and
$\phi_s(U)-t\ge c\eta^2w$ by integration from $r(t)$. Below
$r(t)-w$ the opposite inequality holds, using also the compact
height bound. Relative to the moving root these inequalities give
the same $O(\eps\eta^{-11/2})$ capture bound as above.
An arc already on the inner branch obeys the same tube inequality.
These are times in descending scaled height. In the original flow
$\dd\tau_{\rm phys}/\dd\ell=S/(Lz)$. Throughout the minimum band
$z\asymp1$ for these arcs (the wall bounds $z$ above, and $z=1/2$
points inward from below). The full band therefore takes $O(\eta^{-2})$
in physical time and the forced flight above $r_\mu$ takes $O(\eta^2)$.
The two height increments fit inside the available $2h$, so the arcs share the
landing tube by $\mu-4h$. At the maximum the outer branch
followed through its band is regular, with slope at least $c\eta$
by the simple third root of the limiting cubic.

\emph{Which gap starts the contraction.}
For a launch outside the open window the separation from the reference
at $t_m$ is bounded above and away from zero. If $t_m$ lies inside,
the two arcs follow different branches down to $\mu+4h$.
There the outer root approaches the cubic double root $k$ and the
inner root approaches $-2k$ in $v$, with $k=\sqrt{(2+s)/3}>0$.
Their $U$-separation is comparable to $\sqrt\eta$, since both tube
errors are $O(w)$. Start the secant integral there. This deletes
exactly the disjoint part above the landing in~\eqref{eq:JU-def}
up to the excised band. In either case the logarithm of the starting
separation costs at most $O(\eps|\log\eta|)$.

\emph{Secants and the terminal section.}
On common retained tubes the secant differs from $\phi_s'(r(t))$
by $O(w)$. The secant is bounded during launch, by separation from the
reference when the launched arc is near zero and by a derivative
bound elsewhere. It is bounded on the fold bands as well. On
$8/S\le t\le h$ the trajectories start above the barrier $x=c_{\min}$
and remain above it since $y\ge8$ there. Thus $z\ge c_{\min}/L$,
and even the linear term $O(L^{-5}/z)$ in $\phi_s'$ stays bounded.
At $y=8$ the physical data $x^2/2$ lie in
$[c_{\min}^2/2,2]$. The return to the lower axis for the sextic is
transverse, with a continuous, nonzero derivative on this compact
section times $J$. Compact $C^1$ flow convergence gives upper and
positive lower derivative bounds for the perturbed fields. Hence
$g_s^U$ is the scaled $U$ gap there times $L^2$ and a factor
bounded above and away from zero. The resulting costs in $-\eps\log g_s^U$
are as follows.
\[
\begin{array}{@{}ll@{}}
\toprule
\text{Part of the passage}&\text{Bound on its error}\\ \midrule
\text{Initial separation and terminal conversion}&O(\eps|\log\eta|)\\
\text{Launch approach and common retained tubes}&O(\eta^{9/2}+\eta^{7/2})\\
\text{Two fold bands and the interval }[8/S,h]&O(\eta^3)\\
\text{Omitted action from }C_0/L\text{ to } z_s(h)&
 O(\int_0^{C\sqrt\eta}z^9\,\dd z)=O(\eta^5)\\ \bottomrule
\end{array}
\]
Indeed $\phi_s'(r(t))\,\dd t=(\phi_s'(r))^2\,\dd r$ on each rising
arc. The bound $\phi_s'\le C z^4$ above $C_0/L$ gives the last row.
Integration of~\eqref{eq:secant-exact} now proves
$-\eps\log g_s^U=\J_s^U(t_m)+O(\eta^3)$, with
$\eta^3=o(\eta^{5/2})$.

Finally, write $\Phi_s=\Phi_{\mathrm{ev}}+s\Phi_{\mathrm{od}}$. On $[0,C_0/L]$, substituting $x=Lz\in[0,C_0]$ gives $4\int_0^{C_0/L}(\Phi_{\mathrm{ev}}'(z)/z)\Phi_{\mathrm{od}}'(z)\,\dd z = 4(L^2/S^2)\int_0^{C_0}(F_{\mathrm{ev}}'(x)/x)F_{\mathrm{od}}'(x)\,\dd x = O(\eta^{10})$. On $[0,1]$, expanding $4\int_0^1(\Phi_{\mathrm{ev}}'(z)/z)\Phi_{\mathrm{od}}'(z)\,\dd z$ eliminates orders $\eta,\eta^{3/2},\eta^2$ by~(b), leaves $K(\pi)^2\Lambda\eta^{5/2}$ at order $\eta^{5/2}$, and bounds the $a_4,a_2,a_3,a$ seed cross terms by $O(L^{-2}\eta+L^{-3})=O(\eta^3)$. Adding $\int_{1/2}^{1/2+\sqrt\eta v}(\phi_s')^2\,\dd u=K(\pi)^2\eta^{5/2}P_{-2-s}(v)+o(\eta^{5/2})$—noting that $-2-s$ swaps $(s=+1\leftrightarrow(-3,+8),\ s=-1\leftrightarrow(-1,-8))$ relative to~\eqref{eq:D-cub}—yields $\J_+^U(t_m)-\J_-^U(t_m)=-K(\pi)^2\eta^{5/2}D_{\mathrm{cub}}^U(\xi_m)+o(\eta^{5/2})$ uniformly on $\mathcal{P}$, giving the signs $-D_{\mathrm{cub}}^U(\xi_m)$ in Table~\ref{tab:cubic}.
\end{proof}

\begin{proof}[Proof of Proposition~\ref{prop:low-deg} for $n=10$]
Take $\mathcal{P}$ to be a singleton, write $c=a_6>0$, $b=a_5<0$, $K_{\mathrm{ten}}=-45/4$, $\kappa_S=\frac{8c}{3}$, and $\kappa_L=\frac{b}{\kappa_S K_{\mathrm{ten}}}>0$, and for $0<\eta\ll 1$ put $L=\kappa_L\eta^{-1}$ and $S=\kappa_S L^6$. Define
\begin{align*}
  R(z) &= \tfrac32 z^2(1-z^2), & B(z) &= K_{\mathrm{ten}}z^5(1-z^2)\!\left(1-\tfrac{17}{15}z^2\right), & H(z) &= \tfrac{-594z^7+714z^9}{5},\\
  D_1(z) &= \frac{342z^4-420z^6}{13}, & Z(z) &= \frac{9578943}{2912}z^4, & A_0(z) &= R(z)^2 = \tfrac94 z^4(1-z^2)^2,
\end{align*}
and define $F_{\eta,a}(x)=F_\eta^0(x)+ax$ ($a\in J$) by~\eqref{eq:seeded-template}. Then $\deg F_{\eta,a}=10$ with $[x^{10}]F_{\eta,a}>0$ and $[x^9]F_{\eta,a}<0$. Here $\widetilde Q(z^2)\equiv\frac94=6a_6/\kappa_S>0$ on $[0,z_b]$ for $z_b=5/2$, $E_0(z)=\frac38 z^6(1-\frac32 z^2+\frac35 z^4)$, $\kappa_S\kappa_L[z^5]B=a_5$, and $K=\frac43\cdot\frac94=3$. Direct polynomial integration gives $B(1)=0$, $B'(1)=-3$, $D_1(1)=-6$, $H(1)=24$, $\int_0^1 R^2 B'=\int_0^1 R^2 H'=\int_0^1 D_1 B'=0$, and $4\int_0^1(ZB'+D_1 H')=9\Lambda$. Setting $T_*=E_0(2)+1<E_0(5/2)$, Proposition~\ref{prop:moving-layer} closes $\D(ST_*)=0$ at $a_\eta\in\operatorname{int}J$, retains $+,-,+,-,+$ at $(h_0,\dots,h_4)$, and yields $-,+,-,+,-,+,-,+$ at $St_0<\dots<St_7$. At the regular height $Y=E_0(2)\in(E_0(1),T_*)$, the same $\eta^{7/2}$-tube tracking gives $-\eps\log\big(w_s(ST_*)-w_s(SY)\big)=\J_s^U(Y)+O(\eta^3)$; since each local jump deletion near $z=1$ is $O(\eta^{5/2})$ while the upper root $v_s(Y)=v_e(Y)-s\eta B(2)/(2R(2)^2)+O(\eta^{3/2})$ contributes $-s\eta R(2)^2 B(2)$ to $2s\eta\int_0^2 R^2 B'\,\dd z$, we obtain $\J_+^U(Y)-\J_-^U(Y)=2\eta W_{10}(2)+O(\eta^{3/2})$, where $W_{10}(z)=\int_0^z(R^2 B'-(R^2)'B)\,\dd v=-\frac{9}{16}z^9(z^2-1)^3(3z^2-5)$ satisfies $W_{10}(2)=-54432<0$, giving $\D(SE_0(2))<0$. Ordering $(h_0,\dots,h_4,\,St_0,\dots,St_7,\,SE_0(2),\,ST_*)$ gives $(+,-,+,-,+)+(-,+,-,+,-,+,-,+)+(-)+(0)$; applying Proposition~\ref{prop:sard}(i) with $\sigma_{\mathrm{det}}=+1$ yields $4+1+7+1+1=14$ hyperbolic limit cycles at exact degree $10$.
\end{proof}

\subsection{Closing the decic incoming list in degree thirteen}
The same decic has six strict incoming changes at zero translation.
To transfer the compact incoming ordinates, write $z=x/L$. The ratio
$SE_0(z)/x^6=c(1-3z^2/2+3z^4/5)$ has minimum $c/16>1$ and exceeds
$1+c_0(1+z^4)$ for a fixed $c_0>0$. In this ratio the even controls
cost $O(\eta)(1+z^4)$, the odd controls cost $O(x^{-1})(1+z^4)$
because $\eta L$ is fixed, and the physical lower powers cost $O(x^{-2})$
for $x\ge2$. Choose $X$ large and then use compact convergence on $[2,X]$
to $p_a(sx)>8x^6$. This gives $F_{\eta,a}(sx)\ge x^6$ on the whole
ray $x\ge2$, uniformly in $a\in J$. Truncating at $X$ costs $O(X^{-4})$
in the incoming ordinate. Inward spatial flow contracts ordinate errors,
so Lemma~\ref{lem:incoming-rays} gives convergence at each of $y=1,0$.
For the sextic at zero shift the certified differences satisfy
\[
 -2.4\cdot10^{-9}<I_+(1)-I_-(1)<-1.5\cdot10^{-9},
 \qquad 0.008<I_+(0)-I_-(0)<0.013 .
\]
The $x=2$ enclosures and the same Taylor residual as in
Lemma~\ref{lem:seed-incoming} give these bounds throughout $J$
(\texttt{verify\_entrance.py}). The decic therefore retains both signs.

At the other five ordinates take $Y=E_0(2),t_7,t_6,t_3,t_0$.
Here $t_0$ is the first of the eight cubic samples.
Each level is outside both fold windows. For each side the incoming
gap $g_s^I=I_s(SY)-w_s(ST_*)$ is positive. In scaled coordinates
an incoming orbit from infinity reaches a fixed section to the right of
all reference roots at $Y+O(\eps)$, by the integrable tail
$\eps z/(\Phi_s(z)-t)$ and the uniform decic leading term.
There it is separated from the reference by a fixed positive distance.
It enters the outer rising tube after a further $O(\eps\eta^{-11/2})$
height descent, by the residual bound in Proposition~\ref{prop:moving-layer}.
On the common descent the secant equation, the $4\eta^3$ fold bands and
the conversion at physical height eight in that proof give
$-\eps\log g_s^I=\J_s^I(Y)+O(\eta^3)$. The retained actions at these
five levels agree with the upper actions on each side. At the matched
reference the incoming difference is $g_+^I-g_-^I$, whereas the upper
difference is $g_-^U-g_+^U$. Thus the incoming signs are opposite
to the corresponding upper signs. In decreasing physical ordinate,
\[
 \begin{array}{c|rrrrrrr}
 y& SE_0(2)&St_7&St_6&St_3&St_0&1&0\\ \hline
 \sgn(I_+(y)-I_-(y))&+&-&+&-&+&-&+
 \end{array}
\]
All seven signs persist while detuning and regularizing the decic to
obtain its fourteen compact cycles. Applying Proposition~\ref{prop:cap-transfer}
at a fresh outer scale gives $q=14$, $r=6$, $\sigma=+$, $m=10$, $j=9$,
$b<0$ and $p=1$. The new field has exact degree $\max(12,13)=13$
and at least $14+6+2=22$ simultaneous hyperbolic cycles.

\section{Nineteen cycles in degree twelve}\label{subsec:deg12-prep}
A separate exterior maximum over the decic costs three more degrees by Proposition~\ref{prop:cap-transfer}. To reach nineteen cycles within exact degree twelve, we instead embed the stationary inflection at $z=1$ and the exterior maximum at $z=d$ inside \emph{one} degree-twelve even polynomial and use the geometric radius $d\in D:=[23/20,5/4]$ itself to solve the falling-branch calibration $V(0)=0$.

For $d\in D=[23/20,5/4]$, put $K=\frac43(1-d^{-2})>0$ and define
\begin{equation}\label{eq:deg12-even}
  A(z) = z^4(1-z^2)^2\!\left(1-\frac{z^2}{d^2}\right),\qquad E(z) = \int_0^z vA(v)\,\dd v.
\end{equation}
On $z>0$, $E$ rises through a stationary inflection at $z=1$ ($E(1)>0$) to a nondegenerate global maximum at $z=d$ ($E(d)>E(1)$), and falls strictly on $(d,\infty)$ to $-\infty$; for $Y<E(d)$ let $v(Y)>d$ denote its unique falling root. Write $\M(P)=\int_0^1 A(z)P'(z)\,\dd z$, and for any odd polynomial $B$ put $W(r)=\int_0^r(AB'-A'B)\,\dd z$ and $V(Y)=W(v(Y))-W(d)$. As in the sextic calculation,
$W=-I_1$ for $(f_e,F_o)=(A,B)$ in \cite[Proposition~1]{DMH2015}.
The falling branch retains $[d,v(Y)]$, so its coefficient is
$V(Y)=I_1(d)-I_1(v(Y))$. Set $\omega(z)=z^8(1-z^2)^2(1-z^2/d^2)^2$ and define
\begin{equation}\label{eq:deg12-B0}
  \begin{aligned}
    k &= -\frac{\int_0^1\omega(z)(1+z^2)\,\dd z}{\int_0^1\omega(z)(3z^2-z^4)\,\dd z},\qquad P_k(X)=1+(1+3k)X-kX^2,\\
    B^0(z) &= -z^5(1-z^2)\!\left(1-\frac{z^2}{d^2}\right)\!(1+kz^2).
  \end{aligned}
\end{equation}
On $(0,1)$, $3z^2-z^4>0$ and $(1+z^2)-(3z^2-z^4)=(1-z^2)^2>0$, so $k<-1$. Direct differentiation of $B^0/A= -z(1+kz^2)/(1-z^2)$ gives the exact Wronskian identity
\begin{equation}\label{eq:deg12-wronskian}
  A(z)(B^0)'(z) - A'(z)B^0(z) = -\omega(z) P_k(z^2),
\end{equation}
whence $B^0(1)=B^0(d)=\M(B^0)=0$ and $(B^0)'(1)=2(1-d^{-2})(1+k)<0$. Next put $U(z)=z^7(1-z^2)$ and $U_1(z)=U(z)(1-z^2/d^2)$; since $(U_1/A)'=(z^3/(1-z^2))'>0$ on $(0,1)$, $2\M(U_1)=\int_0^1 A^2(U_1/A)'\,\dd z>0$. Define
\[
  \begin{aligned}
    N_{\mathrm{cap}}(z) &= \frac{U(z)-U_1(z)\M(U)/\M(U_1)}{U(d)}, & \widehat B(z) &= B^0(z) + \theta N_{\mathrm{cap}}(z),\\
    \beta &= \frac{-K}{\widehat B'(1)}, & B(z) &= \beta\widehat B(z).
  \end{aligned}
\]
Then $N_{\mathrm{cap}}(1)=\M(N_{\mathrm{cap}})=0$ and $N_{\mathrm{cap}}(d)=1$. Fix $\theta_0>0$ small so that across the compact rectangle $\mathcal{P}=D\times[-\theta_0,\theta_0]$ we have $\beta(d,\theta)>0$ and
\begin{equation}\label{eq:deg12-B-conds}
  B(1)=\M(B)=0,\qquad B'(1)=-K,\qquad B(d)=\beta\theta,\qquad [z^5]B=-\beta<0.
\end{equation}

\begin{lemma}[Calibration of the exterior maximum radius]\label{lem:deg12-radius}
Put $I_0(d)=\int_d^{v(0)}\omega(z)P_k(z^2)\,\dd z$. At $\theta=0$, $V(0)=-\beta I_0(d)$ has opposite strict signs at the endpoints of $D$, with $I_0(23/20)\in(-466,-464)\cdot 10^{-6}$ and $I_0(5/4)\in(777,780)\cdot 10^{-6}$. Moreover, at every $d\in D$ with $\theta=0$ and $V(0)=0$, one has $V(-1/2)<0$ and $V(Y)>0$ for all $0<Y<E(d)$.
\end{lemma}

\begin{proof}
At $d\in\{23/20,5/4\}$, exact rational integration of the moments in~\eqref{eq:deg12-B0}, $140$-step rational bisection of the strictly decreasing polynomial $E(v)=0$ on $[d,2d]$, and interval evaluation of the polynomial primitive of $\omega(z)P_k(z^2)$ certify the two stated intervals; since $\beta>0$, $V(0)>0$ at $d=23/20$ and $V(0)<0$ at $d=5/4$. Because $k<-1$, the quadratic $P_k(X)=1+(1+3k)X-kX^2$ has $P_k(0)=1>0$, $P_k(1)=2+2k<0$, and leading coefficient $-k>0$, so it has one root in $(0,1)$ and one root $X_*>1$. At any $d\in D$ where $\int_d^{v(0)}\omega(z)P_k(z^2)\,\dd z=0$, positivity of $\omega(z)$ on $(d,v(0))$ forces $d^2<X_*<v(0)^2$. By~\eqref{eq:deg12-wronskian} and $\beta>0$, $AB'-A'B=-\beta\omega(z)P_k(z^2)$ is strictly positive on $(d,\sqrt{X_*})$ and strictly negative on $(\sqrt{X_*},\infty)$. For $0<Y<E(d)$, $d<v(Y)<v(0)$, so either $v(Y)\le\sqrt{X_*}$ (making $V(Y)=\int_d^{v(Y)}(AB'-A'B)>0$ directly) or $\sqrt{X_*}<v(Y)<v(0)$, where $V(Y)=V(0)-\int_{v(Y)}^{v(0)}(AB'-A'B)=-\int_{v(Y)}^{v(0)}(AB'-A'B)>0$. For $Y=-1/2<0$, $v(-1/2)>v(0)>\sqrt{X_*}$, so $V(-1/2)=\int_{v(0)}^{v(-1/2)}(AB'-A'B)<0$.
\end{proof}

Now choose $H(z)=8Kz^7-\mu_H U_1(z)$ with $\mu_H$ fixed by $\M(H)=0$ (so $H(1)=8K$), and set $D_1(z)=-2Kz^4+c_1 B(z)/z$ and $Z(z)=c_2 B(z)/z$, where $c_1,c_2$ are uniquely determined via the positive pivot $\int_0^1(B(z)/z)B'(z)\,\dd z=\frac12\int_0^1 B(z)^2 z^{-2}\,\dd z>0$ by
\begin{equation}\label{eq:deg12-moments}
  \int_0^1 D_1(z)B'(z)\,\dd z = 0,\qquad 4\int_0^1\big(Z(z)B'(z)+D_1(z)H'(z)\big)\,\dd z = K^2\Lambda\qquad\Big(\Lambda=\frac{493}{2}\Big).
\end{equation}
Notice $D_1(1)=-2K$ because $B(1)=0$, while $D_1,Z\in z^4\R[z^2]$, $B\in z^5\R[z^2]$, and $H\in z^7\R[z^2]$ depend smoothly on $\pi=(d,\theta)\in\mathcal{P}$, with $\deg A=10$, $\deg D_1,\deg Z\le10$ and $\deg B=11$, $\deg H\le 11$. With $c=a_6>0$ and $b=a_5<0$ from~\eqref{eq:seed-coeffs}, set $\kappa_S=6c$ and $\kappa_L(d,\theta)=\frac{b}{6c(-\beta(d,\theta))}>0$, and couple the scales for $0<\eta\ll 1$ by
\begin{equation}\label{eq:deg12-poly}
  \begin{aligned}
    L &= \kappa_L(d,\theta)\eta^{-1} > 0,\qquad S = 6c L^6,\qquad \eps = \frac{L^2}{S^2}\asymp\eta^{10},\\
    E_\eta(z) &= E(z) + \int_0^z v\big(\eta D_1(v)+\eta^{3/2}Z(v)\big)\,\dd v,\\
    F_{\eta,a}(x;d,\theta) &= S\big(E_\eta(x/L)+\eta B(x/L)+\eta^{3/2}H(x/L)\big) + a_2 x^2 + a_3 x^3 + a_4 x^4 + ax,
  \end{aligned}
\end{equation}
which has exact degree twelve ($[x^{12}]F_{\eta,a}=SL^{-12}(-1/(12d^2)+O(\eta))<0$) and fits~\eqref{eq:seeded-template} identically with $E_0=E$, $A_0=A$, and $\widetilde Q(z^2;d)=1-z^2/d^2$.

Put $T(d)=(E(1)+E(d))/2$ and define $z_b(d)\in(1,d)$ continuously on $D$ by $E(z_b(d))=(T(d)+E(d))/2$. Since $\max_{d\in D} z_b(d)/d<1$, we have $\min_{(d,\theta)\in\mathcal{P},\,z\in[0,z_b(d)]}\widetilde Q(z^2;d)=\min_{d\in D}(1-z_b(d)^2/d^2)>0$, while~\eqref{eq:deg12-B-conds} and~\eqref{eq:deg12-moments} verify all hypotheses~(a)--(b) of Proposition~\ref{prop:moving-layer} on $\mathcal{P}$. Consequently, for small $\eta>0$ there is a unique jointly continuous $a=a_\eta(d,\theta)\in\operatorname{int}J$ closing $\D(ST(d))=0$ ($w_*:=w_s(ST(d))$), at which $(h_0,\dots,h_4)$ have signs $+,-,+,-,+$ and $St_0<\dots<St_7$ have signs $-,+,-,+,-,+,-,+$.

In normalized variables $z=x/L$, $u=z^2/2$, $t=y/S$, the backward
equation is $\eps\,\dd u/\dd t=\phi_s(u)-t$, with
$\phi_s(z^2/2)=\Phi_s(z)=F_{\eta,a}(sLz;d,\theta)/S$.
The orbit from $(u,t)=(0,-1)$ hits the upper axis at scaled height
$H_s^*$. An orbit from the section $z_e=1/2$ at physical ordinate $y$
hits at $h_s(y)$, with gap $g_s(y)=H_s^*-h_s(y)>0$. The launch from
$(0,-1/2)$ gives the analogous gap $g_s^b$.

For the limiting profile, $J_o(Y)=\int_d^{v(Y)}zA(z)^2\,\dd z$
and $J_o'(Y)=A(v(Y))<0$. Write $P_e=(\phi_++\phi_-)/2$ and
$\Phi_{\mathrm{ev}}(z)=P_e(z^2/2)$ for the actual even profile in the
two coordinates. Its maximum near $d$ has spatial coordinate
$\widehat d_e$ and height $M_e=\Phi_{\mathrm{ev}}(\widehat d_e)$.
Let $v_e(Y)>\widehat d_e$ be the falling root
$\Phi_{\mathrm{ev}}(v_e(Y))=Y$. Since $\dd u=z\,\dd z$ and
$P_e'(z^2/2)=\Phi_{\mathrm{ev}}'(z)/z$, the positive falling action is
\[
 J_e(Y)=\int_{\widehat d_e}^{v_e(Y)}
           \frac{\Phi_{\mathrm{ev}}'(z)^2}{z}\,\dd z
       =\int_{\widehat d_e^{\,2}/2}^{v_e(Y)^2/2} P_e'(u)^2\,\dd u .
\]

The regular approach at a fixed section $b_0\in(0,E(z_e))$
gives~\eqref{eq:entry-distortion} and a bundle of width
$O(e^{-c_1/\eps})$, uniformly by compactness. Near $d$ the maximum
has curvature bounded away from zero. A tube of radius $\eta^3$
tracks the falling root to $4\eta^2$ below its actual maximum,
since $|\phi_s'|\ge c_6\eta$, the root speed is $O(\eta^{-1})$
and $\eta^4/\eps\gg\eta^{-1}$. Above the maximum plus $\eta^2$,
the remaining flight takes $O(\eps/\eta^2)$ in height.
Lemma~\ref{lem:terminal-distortion} thus applies with $D=O(\eta^2)$
and a fixed inverse strip near the upper-axis hit. At physical
resolution the entry error costs
$O(S\eps|\log\eps|)=O(\eta^4|\log\eta|)=o(1)$.
In $J_e$ the integrand vanishes at the maximum. At the regular falling
root the endpoint displacement contributes $-s\eta A(v(Y))B(v(Y))$
to the cross term $2s\eta\int_d^{v(Y)}AB'\,\dd z$. Their sum is
$s\eta V(Y)$ in the positive action, hence $-s\eta V(Y)$ in the log gap.
Consequently, uniformly on $\mathcal P$, $a\in J$, and bounded $y,y_0$,
\begin{align}
  H_s^* &= M_e + s\eta\beta\theta + O(\eta^{3/2}),\qquad
  \eps\log g_s(y) = -J_e(0) - s\eta V(0) + O(\eta^{3/2}),\label{eq:deg12-cap}\\
  \eps\log g_s^b &= -J_e(-1/2) - s\eta V(-1/2) + O(\eta^{3/2}),\qquad
  S\eps\log\frac{g_s(y)}{g_s(y_0)} = -A(v(0))(y-y_0) + o(1).\label{eq:deg12-fine}
\end{align}
Substituting $a=a_\eta(d,\theta)$ on the rectangle $D\times[-\eta^{1/4},\eta^{1/4}]$ ($\eta^{5/4}\gg\eta^{3/2}$), $H_+^*-H_-^*$ has opposite strict signs on the $\theta=\pm\eta^{1/4}$ faces by~\eqref{eq:deg12-cap}, while $\eps\log(g_+(\Delta)/g_-(0))$ ($\Delta=3/500$) has opposite strict signs on the $d=23/20$ and $d=5/4$ faces by Lemma~\ref{lem:deg12-radius}. The two discrepancies are continuous in $(d,\theta)$. By the rectangle form of Brouwer's fixed-point theorem, these face signs give an interior solution $(d_\eta,\theta_\eta)$ satisfying simultaneously
\begin{equation}\label{eq:deg12-3matches}
  \D(ST)=0,\qquad H_+^*=H_-^*=:H_*,\qquad g_+(\Delta)=g_-(0),
\end{equation}
along which $\theta_\eta\to 0$ and $V(0)\to 0$ as $\eta\downarrow 0$. Just as in Proposition~\ref{prop:cap-transfer},~\eqref{eq:deg12-fine}, Lemma~\ref{lem:incoming-rays}(ii) on $[2,Lz_e]$, and Lemma~\ref{lem:seed-incoming} yield upper-axis signs $-,+,-$ at $S h_+(y)$ for $y=4/5,0,-20$ in increasing upper height, followed before $SH_*$ by a direct lower-axis sample from $t=-1/2$ of strict sign $-\sgn V(-1/2)=+$ (since $V(-1/2)<0$ uniformly near the compact limiting match set $\{\theta=0,V(0)=0\}$ by Lemma~\ref{lem:deg12-radius}).

\subsection{The turning orbit between the matched references}\label{subsec:deg12-finish}
\begin{figure}[htbp]
\centering
\begin{tikzpicture}[x=1.17cm,y=0.94cm,>=Stealth,font=\small]
\draw[->] (0,0)--(5.05,0) node[right] {$z$};
\draw[->] (0,0)--(0,2.75) node[left] {$t$};
\draw[thick] (0,0) .. controls (.38,0) and (.59,.12) .. (.84,.51)
 .. controls (1.26,1.2) and (1.64,1.33) .. (2.02,1.33)
 .. controls (2.55,1.33) and (2.85,2.3) .. (3.35,2.3)
 .. controls (3.9,2.3) and (4.12,1.30) .. (4.38,.51)
 .. controls (4.46,.25) and (4.5,.12) .. (4.56,-.16);
\draw[densely dotted] (0,.51) node[left] {$Y$}--(.84,.51);
\draw[very thick,dashed,->] (4.38,.51)--(.91,.51);
\draw[very thick,->] (.84,.51) .. controls (.72,.32) and (.57,.16) .. (.43,.105);
\fill (.84,.51) circle (1.5pt) (4.38,.51) circle (1.5pt);
\draw (2.02,.04)--(2.02,-.04) node[below] {$1$};
\draw (3.35,.04)--(3.35,-.04) node[below] {$d$};
\node[below] at (.84,0) {$r$}; \draw (4.38,.04)--(4.38,-.04) node[below left,inner sep=2pt] {$v(Y)$};
\node[above] at (3.35,2.3) {$E(d)$}; \node[above] at (2.02,1.33) {$E(1)$};
\end{tikzpicture}
\caption{Turning launch at $Y=E(r)$, shown on the limiting even profile
(schematic). In forward time in the reflected right half-plane the
trajectory leaves the falling root $v(Y)>d$, moves inward to the rising
root $r<1$ at this limiting height, and follows the inner branch toward
the axis. The backward half (omitted) follows the falling branch in
increasing $t$ and reaches the upper axis through the distant maximum.}\label{fig:turning-supp}
\end{figure}

At~\eqref{eq:deg12-3matches}, the last inflection sign ($+$) and the first section-entrance sign ($-$) frame the matched inner reference $\D(ST)=0$. If we detune $a$ upward so that $\D(SH_*)<0$ (sign $-$ after the $+$ from $t=-1/2$), then $\D(ST)$ also becomes negative. The first section-entrance sign is also negative. Two more brackets
therefore require a positive sign between $ST$ and $S h_+(4/5)$. A regular turning orbit on the outer falling branch supplies that positive bridge.

For a small fixed $r\in(0,1/2)$, put $Y=E(r)\in(0,E(1))$ and, on each side $s\in\{\pm 1\}$ at a match of~\eqref{eq:deg12-3matches}, launch an orbit at its outer turning point $(z,t)=(v_s(Y),Y)$ where $\Phi_s(v_s(Y))=Y$ on the falling branch ($v_s(Y)>d$). Denote its backward scaled upper hit by $H_s(Y)$ and its forward physical lower-axis hit magnitude by $b_s(Y)$. Let $I_e(Y)=\int_{C_0/L}^{r_e(Y)}\Phi_{\mathrm{ev}}'(z)^2\,\dd z/z$ be the rising action in the $z$ coordinate, ending at the regular inner root $\Phi_{\mathrm{ev}}(r_e(Y))=Y$ near $r$. Uniformly for levels in a small neighborhood of $Y$, we claim that $H_s(Y)<H_*$, $b_s(Y)>w_*$, and
\begin{align}
  \eps\log\big(H_* - H_s(Y)\big) &= -J_e(Y) - s\eta V(Y) + O_r(\eta^{3/2}),\label{eq:turn-back}\\
  \eps\log\big(b_s(Y) - w_*\big)  &= -I_e(Y) - s\eta W(r) + O_r(\eta^{3/2}).\label{eq:turn-fwd}
\end{align}

To prove~\eqref{eq:turn-back}, run~\eqref{eq:scalar-u} in backward time ($t$ increasing). The outer reference $u_{\mathrm{ref}}^{\mathrm{out}}(t)$ from $(u,t)=(0,-1)$ tracks the falling root $u_s(t)=v_s(t)^2/2$ from lower heights ($t<Y$) while remaining strictly \emph{above} $u_s(t)$ in $u$ (to the right of the falling graph in $(z,t)$): indeed, $e(t)=u_{\mathrm{ref}}^{\mathrm{out}}(t)-u_s(t)>0$ satisfies $e'(t)=[\phi_s]_{u_{\mathrm{ref}}^{\mathrm{out}},u_s}\,e(t)/\eps - u_s'(t)$, where both $-[\phi_s]_{u_{\mathrm{ref}}^{\mathrm{out}},u_s}$ and $-u_s'(t)$ are bounded above and away from zero on a compact regular approach interval $[Y-\delta_r,Y]$. Variation of constants after entry into the falling-branch tube therefore places $e(Y)=u_{\mathrm{ref}}^{\mathrm{out}}(Y)-u_s(Y)$ strictly between $c_r\eps$ and $C_r\eps$ ($0<c_r<C_r$). Since the turning orbit starts at $u_{\mathrm{turn}}(Y)=u_s(Y)<u_{\mathrm{ref}}^{\mathrm{out}}(Y)$, trajectory order preserves this positive gap to the upper axis ($H_s(Y)<H_s^*=H_*$), and since $\eps|\log\eps|=O(\eta^{10}|\log\eta|)=o(\eta^{3/2})$, applying the upper-hit derivative~\eqref{eq:Ks-deriv} and falling-action expansion of~\eqref{eq:deg12-cap} from $t=Y$ yields~\eqref{eq:turn-back}.

To prove~\eqref{eq:turn-fwd}, follow $u_{\mathrm{turn}}(t)$ and the inner reference $u_{\mathrm{ref}}^{\mathrm{in}}(t)$ from $(0,ST)$ in forward time ($t$ decreasing from $Y$). Because $Y=E(r)<E(1)$ lies strictly below the stationary inflection and $T>E(1)>Y$, at the launching ordinate $t=Y$ the inner reference is already tracking the inner rising root $\tilde u_s(t)=r_{s,\eta}(t)^2/2$ near $r^2/2$, whereas $u_{\mathrm{turn}}(Y)=u_s(Y)=v_s(Y)^2/2$ starts near $v(Y)^2/2>d^2/2>r^2/2$. Thus, \emph{at the launching ordinate $t=Y$ itself}, $u_{\mathrm{turn}}(Y)-u_{\mathrm{ref}}^{\mathrm{in}}(Y)=\frac12(v(Y)^2-r^2)+o(1)\ge c_r'>0$ is bounded above and away from zero; trajectory order preserves $u_{\mathrm{turn}}(t)>u_{\mathrm{ref}}^{\mathrm{in}}(t)$ down to $\Sigma_-$, proving $b_s(Y)>w_s(ST)=w_*$ and bounding the initial term $\eps\log\big(u_{\mathrm{turn}}(Y)-u_{\mathrm{ref}}^{\mathrm{in}}(Y)\big)=O_r(\eps)$ in~\eqref{eq:secant-exact}. Next, since $\Phi_s(z)>Y$ across the open interval $(r_{s,\eta}(Y),v_s(Y))$, fix $\nu_r>0$ small so that $-\phi_s'\ge 2c_r>0$ and $-u_s'\ge 2c_r>0$ on $[u_s(t)-\nu_r,u_s(t)]$, $\phi_s'\ge c_r''>0$ and $|\tilde u_s'|\le C_r''$ on $[\tilde u_s(t),\tilde u_s(t)+\nu_r]$, and $\phi_s\ge Y+\kappa_r$ ($\kappa_r>0$) across $[\tilde u_s(Y)+\nu_r,u_s(Y)-\nu_r]$:
\begin{enumerate}
  \item[\textup{1.}] \emph{Falling-root escape:} while $\tilde e(t):=u_s(t)-u_{\mathrm{turn}}(t)\in[0,\nu_r]$, $\dd\tilde e/\dd(-t)=-u_s'(t)-[\phi_s]_{u_{\mathrm{turn}},u_s}\tilde e/\eps\ge c_r+c_r\tilde e/\eps$ from $\tilde e(Y)=0$, reaching $\tilde e=\nu_r$ in height descent $\tau_1\le(\eps/c_r)\log(1+\nu_r/\eps)=O_r(\eps|\log\eps|)$;
  \item[\textup{2.}] \emph{Horizontal bridge:} across $[\tilde u_s(Y)+\nu_r,u_s(Y)-\nu_r]$, $\dd u_{\mathrm{turn}}/\dd(-t)=(t-\phi_s(u_{\mathrm{turn}}))/\eps\le -\tfrac12\kappa_r/\eps$, taking height descent $\tau_2=O_r(\eps)$;
  \item[\textup{3.}] \emph{Inner rising-root capture:} on $[\tilde u_s(t),\tilde u_s(t)+\nu_r]$, $q(t):=u_{\mathrm{turn}}(t)-\tilde u_s(t)>0$ obeys $\dd q/\dd(-t)\le C_r''-c_r'' q/\eps$, entering an $O_r(\eps)$ neighborhood (and hence the $\eta^{7/2}$-tube of Proposition~\ref{prop:moving-layer}) within $\tau_3=O_r(\eps|\log\eps|)$ height descent.
\end{enumerate}
Throughout the entire transfer interval $[Y-(\tau_1+\tau_2+\tau_3),Y]$ of height length $O_r(\eps|\log\eps|)$, both $u_{\mathrm{turn}}$ and $u_{\mathrm{ref}}^{\mathrm{in}}$ remain inside the compact positive interval $[\frac14 r^2,v(0)^2]$ where $\phi_s'$ and $[\phi_s]_{u_{\mathrm{turn}},u_{\mathrm{ref}}^{\mathrm{in}}}$ are $O_r(1)$, so~\eqref{eq:secant-exact} changes by $O_r(\eps|\log\eps|)=o(\eta^{3/2})$. Below $Y-(\tau_1+\tau_2+\tau_3)$, Proposition~\ref{prop:moving-layer} applies with \emph{no} intervening fold on $[C_0/L,r_{s,\eta}(Y)]$ (all compact seed critical points lie below $z=C_0/L$ at height $O(\eta^6)\ll\eta^3$), giving $\eps\log(b_s(Y)-w_*)=-\int_{C_0/L}^{r_{s,\eta}(Y)}\Phi_s'(z)^2\,\dd z/z+O_r(\eta^3)$. Expanding the upper integration limit $r_{s,\eta}(Y)=r_e(Y)-s\eta B(r)/E'(r)+O_r(\eta^{3/2})$ contributes the moving-endpoint term
\[
  \big(r_{s,\eta}(Y)-r_e(Y)\big)\frac{E'(r)^2}{r} + O_r(\eta^{3/2}) = -s\eta\,A(r)B(r) + O_r(\eta^{3/2}),
\]
which combines with the cross term $2s\eta\int_0^r A(z)B'(z)\,\dd z+O_r(\eta^{3/2})$ (the $[0,C_0/L]$ cutoff costing $O(\eta^{10})$ and $(\Phi_{\mathrm{od}}')^2/z$ costing $O_r(\eta^2)$) to give $s\eta W(r)=-s\eta I_1(r)$ in the positive rising action.
It enters~\eqref{eq:turn-fwd} with the opposite sign.

Because an earlier falling-branch trajectory already lies strictly to the right when a later turning orbit launches, $H_s(Y)$ decreases strictly with $Y$. At $Y=E(r)$, differentiating the actions at their regular spatial endpoints gives
\[
 J_e'(Y)=\frac{\Phi_{\mathrm{ev}}'(v_e(Y))}{v_e(Y)}
        =A(v(Y))+O_r(\eta)<0,\qquad
 I_e'(Y)=\frac{\Phi_{\mathrm{ev}}'(r_e(Y))}{r_e(Y)}
        =A(r)+O_r(\eta).
\]
The first derivative stays negative near $Y$, so~\eqref{eq:turn-back}
brackets a unique $Y_-=Y+O_r(\eta)$ with $H_-(Y_-)=H_+(Y)$.
Expanding~\eqref{eq:turn-back} and then subtracting~\eqref{eq:turn-fwd} yields
\begin{align}
  Y_- - Y &= 2\eta\,\frac{V(Y)}{A(v(Y))} + O_r(\eta^{3/2}),\label{eq:shift-turn}\\
  \eps\log\frac{b_+(Y)-w_*}{b_-(Y_-)-w_*} &= -2\eta\!\left(W(r) - \frac{A(r)V(Y)}{A(v(Y))}\right) + O_r(\eta^{3/2}).\label{eq:bridge}
\end{align}
On the compact limiting match set $\{\theta=0,\ V(0)=0,\ d\in D\}$, $-A(v(0))$ and $\beta$ have positive minima, and as $r\downarrow 0$ we have $A(r)=r^4+O(r^6)$ and $B(r)=-\beta r^5+O(r^7)\implies W(r)=-\frac19\beta r^9+O(r^{11})$, while smoothness of the regular falling root $v(\cdot)$ near $0$ with $V(0)=0$ gives $V(E(r))=O(E(r))=O(r^6)$ and hence $A(r)V(E(r))/A(v(E(r)))=O(r^{10})$. Fixing $r>0$ sufficiently small first (before choosing $\eta>0$ small) makes $W(r)-A(r)V(Y)/A(v(Y))<0$ uniformly near the limiting match set, so~\eqref{eq:bridge} gives $b_+(Y)>b_-(Y_-)$, i.e., $\D(SH_+(Y))>0$.

\begin{remark}[Explicit rational turning radius]\label{rem:rational-turn}
At $\theta=0$ on the narrowed rational interval $D'=[1.19888999543,\,1.19888999544]\subset D$, $V(0)$ already has opposite strict signs at the two endpoints, and at the explicit rational radius $r=1/4$, monotone interval arithmetic certifies $\beta^{-1}\big(W(1/4)-A(1/4)V(E(1/4))/A(v(E(1/4)))\big)\in(-3.1\cdot 10^{-7},-2.8\cdot 10^{-7})<0$ uniformly across all $d\in D'$. The strict face and bridge inequalities persist at finite matches with $\theta_\eta\to0$ when the matching rectangle uses $D'$.
\end{remark}

\begin{proof}[Completion of the proof of Theorem~\ref{thm:deg12}]
Fix a sufficiently small $\eta>0$ and a matched solution $(d_\eta,\theta_\eta)$ of~\eqref{eq:deg12-3matches} with $a_*=a_\eta(d_\eta,\theta_\eta)\in\operatorname{int}J$, and put $h_{\mathrm{dir}}=H_*-g_-^b$. Consider the $20$ physical upper heights $(H_1,\dots,H_{20})$ defined by
\[
  \big(h_0,\dots,h_4,\ St_0,\dots,St_7,\ ST(d_\eta),\ SH_+(Y),\ S h_+(4/5),\ S h_+(0),\ S h_+(-20),\ Sh_{\mathrm{dir}},\ SH_*\big).
\]
The compact heights end at $h_4=4.36$, while
$St_0=S(E(1)+o(1))\to\infty$. Since $K>0$, the order of the eight
$St_i$ is the order of the cubic samples $\xi_i$. The next two
inequalities follow from
$T(d_\eta)-t_7=(E(d_\eta)-E(1))/2+o(1)>0$
and $H_+(Y)=E(d_\eta)+o(1)>T(d_\eta)$.
For the remaining heights,
$d_\eta<v(Y)<v(0)<v(-1/2)$ gives
$0<J_o(Y)<J_o(0)<J_o(-1/2)$. The absolute actions
in~\eqref{eq:deg12-cap} and~\eqref{eq:turn-back} separate the turning,
entrance and direct heights. Within the entrance list,
$-A(v(0))>0$ in~\eqref{eq:deg12-fine} orders the three gaps. Together,
\[
 H_*-H_+(Y)\gg g_+(4/5)>g_+(0)>g_+(-20)\gg g_-^b>0 .
\]
This orders all twenty heights. At $a=a_*$ their mismatches are
\[
\begin{array}{c|ccccccc}
k&1\text{--}5&6\text{--}13&14&15&16\text{--}18&19&20\\ \hline
\sgn\D(H_k,a_*)&{+}{-}{+}{-}{+}&{-}{+}{-}{+}{-}{+}{-}{+}&0&+&{-}{+}{-}&+&0 .
\end{array}
\]
Now hold $(d_\eta,\theta_\eta,\eta)$ (hence $a_2,\dots,a_{12}$) and \emph{all twenty physical heights $H_1,\dots,H_{20}$ fixed}, and vary only the linear coefficient $a$ upward on a small open interval $(a_*,a_*+\delta_0)\subset\operatorname{int}J$. Since $|F_{\eta,a_*}'(0)|=|a_*|<2$ and $(0,H_{20})=(0,SH_*)$ has transverse lower-axis hits $(0,-S)$ on both sides, Lemma~\ref{lem:half-returns}(iii) gives smooth transverse half-returns $w_\pm(h,a)$ with $w_\pm'>0$ across all of $(0,H_{20}]$ for $|a-a_*|<\delta_0$. By Lemma~\ref{lem:rotational}, $\partial_a\D(h,a)<0$, so for every $a\in(a_*,a_*+\delta_0)$ both zeros at $H_{14}=ST(d_\eta)$ and $H_{20}=SH_*$ become strictly negative while all $18$ nonzero strict signs persist, yielding $20$ strictly alternating signs $+,-,+,-,\dots,+,-$ and $19$ pairwise disjoint open brackets $(H_{k-1},H_k)$ for $2\le k\le20$. Applying Proposition~\ref{prop:sard}(i) ($\sigma_{\mathrm{det}}=-1$) selects a Sard regular value $a_{\mathrm{reg}}\in(a_*,a_*+\delta_0)$ at which every bracket contains a simple zero $r_k$ ($\partial_h\D(r_k,a_{\mathrm{reg}})\neq 0$, multiplier $w_+'(r_k)/w_-'(r_k)\neq 1$), giving $19$ distinct hyperbolic limit cycles on a nonempty open exact-degree-twelve neighborhood containing $\Q^{12}$ points.
\end{proof}

\subsection{An unsplit inner profile in degree nine}\label{subsec:deg9}
\begin{proof}[Completion of Proposition~\ref{prop:low-deg} for $n=9$]
Write $c=a_6>0$, $b=a_5<0$ and fix an outer radius $L>2$. Use
\[
 \begin{aligned}
 A(z)&=6c z^4(1-z^2/L^2),& E(z)&=\int_0^z v A(v)\,\dd v,\\
 B_{\lambda,\theta}(z)&=b z^5(1-z^2/L^2)
     +\lambda L^5(z/L)^7(1-z^2/L^2)+\theta L^5(z/L)^7 .
 \end{aligned}
\]
The maximum of $E$ is at $L$. For $Y<E(L)$ let $v(Y)>L$ be the
falling root and use the same positive-action convention as in
Section~\ref{subsec:deg12-prep}, with $W(r)=\int_0^r(AB'-A'B)\,\dd z$,
$V(Y)=W(v(Y))-W(L)$. Set $z=LZ$. Dividing $A,E,B$ by
$L^4,L^6,L^5$ gives the profiles of Proposition~\ref{prop:cap-transfer}
for $m=6,j=5$, with $v(0)/L=\sqrt{4/3}$. At $\theta=0$ the
normalized density of $V$ has the sign of $b+3\lambda Z^2$ on
$Z>1$. Its integral at $Y=0$ is affine in $\lambda$, with positive
slope and negative value at $\lambda=0$. Thus $V_{\lambda_0,0}(0)=0$
at a unique $\lambda_0>0$, independent of $L$. A bounded $Y$ displaces
the normalized root by $O(|Y|L^{-6})$. Since $V$ scales as $L^9$ and
$|A(v(Y))|\asymp L^4$, we have
$|V_{\lambda_0,0}(Y)/A(v(Y))|\le C|Y|/L$.

On the inner branch take $r=1$, $Y=E(1)$. As $L\to\infty$,
$W(1)=\frac23 cb+O(L^{-2})$ and
$W(1)-A(1)V(Y)/A(v(Y))=\frac23 cb+O(L^{-1})$ at this calibration.
Fix $L$ large enough that both are strictly negative.
Choose $T=(Y+E(L))/2$ and a section $1<z_e<L$ with $E(z_e)>T$.
For $\delta=R^{-1}$, $\eps=\delta^{10}$, set
\[
 F_{R,a}(x)=R^6 E(x/R)+R^5 B_{\lambda,\theta}(x/R)
                  +a_4x^4+a_3x^3+a_2x^2+ax .
\]
This converges to $p_a$ on fixed physical intervals. On $0<z\le z_e$,
$E(z)/z^6$ has positive infimum and $z|B(z)|/E(z)$ is bounded
uniformly for bounded parameters. These bounds give $F_{R,a}(sx)\ge c_1x^6$
on $X\le x\le Rz_e$ after $X$ and then $R$ are large.
The wall at $Rz_e$ traps the upper reference $R^6T$.
Capture at $y=8$ and the two faces in Lemma~\ref{lem:seed} close that
reference at a unique continuous $a_R(\lambda,\theta)\in\operatorname{int}J$,
by the strict derivative~\eqref{eq:rot-half}.

With $L$ fixed, the coarse and fine falling-branch estimates
(3.15) and (3.17) of the main paper apply in the $R$ scaling.
Indeed the scaled even remainder is $O(\delta^2)$, the odd remainder
after $\delta B$ is $O(\delta^3)$ on fixed intervals, and
$\eps|\log\eps|=o(\delta^6)$. The peak splitting has coefficient
$B(L)=\theta L^5$, and $\partial_\lambda V(0)>0$.
Consequently the two residuals have opposite signs on the respective
faces of $[\lambda_0-\rho_0,\lambda_0+\rho_0]\times
[-\sqrt\delta,\sqrt\delta]$, with $0<\rho_0<\lambda_0/2$ fixed.
The rectangle argument of Proposition~\ref{prop:cap-transfer} gives
$H_+^*=H_-^*=H_*$ and $g_+(\delta^6\Delta)=g_-(0)$ for
$\Delta=3/500$, at $\theta_R=O(\delta)$ and
$\lambda_R=\lambda_0+O(\delta)>0$. In particular
$[x^9]F_{R,a}=-\lambda_R R^{-4}L^{-4}<0$, and every higher power vanishes.

The upper mismatch $\D(R^6Y)$ has sign $W(1)<0$ by the regular rising-root
calculation of Section~\ref{sec:deg12}. For the turning launch at
$(v_s(Y),Y)$, equality of the backward upper hits shifts the minus
level by $2\delta V(Y)/A(v(Y))+O(\delta^2)$.
The proof of~\eqref{eq:bridge} applies at these regular roots: the
forward passage has no folds before $r=1$, the backward passage has one
outer maximum, and the same two matched references bound both gaps.
The resulting mismatch has sign
$-\sgn(W(1)-A(1)V(Y)/A(v(Y)))=+$. Thus the samples at $R^6Y$ and
$R^6H_+(Y)$ give successive negative and positive signs. The latter
height exceeds $R^6T$ since $H_+(Y)\to E(L)>T$.
Since $J_o(Y)<J_o(0)$, all three paired section heights are higher
still and below $R^6H_*$. They occur in the order $4/5,0,-20$ by the positive slope in (3.17) of the main paper. The uniform bound $F_{R,a}(sx)\ge c_1x^6$
on $[X,Rz_e]$ gives a tail error $O(X^{-4})$ and transfers all three
strict signs of Lemma~\ref{lem:seed-incoming}. The five compact
signs persist by compact convergence. Holding all eleven upper
heights fixed gives in increasing order
\[
  (+,-,+,-,+),\quad -,\quad +,\quad (-,+,-),\quad 0 .
\]
The final zero is the upper hit from the common lower point $(0,-R^6)$.
Decrease $a$ slightly in $J$ to make it positive and apply
Proposition~\ref{prop:sard}(i) on the trapped return interval.
The ten brackets give ten simultaneous hyperbolic cycles on an open
set of exact degree nine.
\end{proof}

\begin{remark}[Degree-twelve ansatz constraints]\label{rem:deg12-obstruction}
Within the ansatz of Section~\ref{subsec:deg12-prep} ($A$ as in~\eqref{eq:deg12-even} and $B$ odd with quintic germ satisfying $B(1)=B(d)=0$), any polynomial $B$ of degree at most $9$ is a scalar multiple of $z^5(1-z^2)(1-z^2/d^2)$, whose Wronskian with $A$ equals a nonzero multiple of $\omega(z)(1+z^2)$ and therefore cannot satisfy $\M(B)=0$; degree $11$ in $B$ is minimal for that ansatz. Likewise, Lemma~\ref{lem:deg12-radius} gives $V(E(1))>0$ at every limiting match, so~\eqref{eq:shift-turn} shifts the cross-side turning height near $E(1)$ by $O(\eta)$, strictly larger than the $O(\eta^{3/2})$ width of the split cubic windows, preventing direct transfer of the four incoming cubic changes inside this joint degree-twelve family.
\end{remark}
